\chapter{General Equilibrium}

The unifying idea across all market analyses is \emph{market clearing}.
What distinguishes the two standard approaches is how many markets
they ask to clear at once. \emph{Partial equilibrium} fixes prices in
all other markets and analyzes a single market in isolation.
\emph{General equilibrium} requires every market in the economy to
clear simultaneously, taking into account the cross-market price
linkages.

Partial equilibrium remains a useful shortcut despite its narrower
scope. It is a defensible approximation whenever:
\begin{itemize}
  \item
    Prices in all other markets are essentially unaffected by what
    happens in the market of interest (e.g., the market is small
    relative to the rest of the economy).
  \item
    There are no significant wealth effects feeding back from the
    market of interest into demand elsewhere.
  \end{itemize}

The following example shows what can go wrong when these conditions
fail.

\ex{
There are \(n\) towns in total, with $n$ being a large number. Each town has an identical price-taking firm that produces a single consumption good with the
production function \(y = f\pr{l}\), where \(f'\pr{l}>0\) and
\(f''\pr{l}<0\). The single consumption good is traded in the national
market. Suppose there are \(L\) units of inelastic labor supply in total. Workers
can move freely across towns to seek the highest possible wage, which are regarded as complete
information. Normalize \(p_c = 1\) and let \(w_i\) denote the
equilibrium wage in town \(i\). Now suppose that town 1 levies a \textbf{small}
tax \(t > 0\) on firm 1 for each unit of labor hired. Which group of
individual(s) will bear the tax burden?
}

\sol{
\begin{itemize}
\item
  Equilibrium without tax

  \begin{itemize}
  \item
    Whether we adopt the partial or general equilibrium approach, the
    competitive equilibrium without the labor tax is identical.
  \item
    Since workers can move freely and \(f'\pr{l}>0\) and
    \(f''\pr{l}<0\), we must have

    \[\bc
    l^*  = \frac1n \cdot L\\
    w_1  = w_2 =  \cdots  = w_n = f'\pr{l^*} = f'\pr{\frac{l}{n}}
    \ec\]
  \end{itemize}
\item
  Partial equilibrium in town 1 with tax

  \begin{itemize}
  \item
    Assuming the wage rate in other markets are unaffected, we must have
    \(w_1\pr{t} = w_0+t\).
  \item
    At the new equilibrium, \(f'\pr{l_1\pr t} = w_1\pr t = w_0+t\).
  \item
    Hence we can see that firm 1 bears all the tax burden.
  \end{itemize}
\item
  General equilibrium with town 1 levied with tax

  \begin{itemize}
  \item
    Let \(l_1\pr t\) denote the equilibrium amount of labor in town 1
    with a \(t\) unit tax, and \(l_{-1}\pr t\) denote the equilibrium
    amount of labor in any other town when town 1 imposes a \(t\) unit
    tax and \(w\pr t\) the equilibrium wage rate received by the worker.
    When all labor markets clear, we should have:
    \[\ba
    & \bc
    l_1\pr t+\pr{n-1}l_{-1}\pr t=  L\\
    f'\pr{l_1\pr t} = w_1\pr t = w\pr t+t\\
    f'\pr{l_{-1}\pr t} = w\pr t
    \ec\\
    \too & \bc
    f'\pr{L-\pr{n-1}l_{-1}\pr t} = w\pr t+t\\
    f'\pr{l_{-1}\pr t} = w\pr t
    \ec\\
    \too & \bc
    -\pr{n-1}l'_{-1}\pr t\cdot f''\pr{L-\pr{n-1}l_{-1}\pr t} = w'\pr t+1\\
    l'_{-1}\pr tf''\pr{l_{-1}\pr t} = w'\pr t
    \ec
    \ea\]
  \item
    Set \(t\to 0\), so we have \(l_{-1}\pr 0 = \frac1n\cdot L\). From
    the two equations above we obtain \(w'\pr{0} = -\frac1n\).
  \item
    Let \(\Pi\pr{w\pr t}\) denote the total profit of the firms when the
    wage rate is \(w\pr t\) and \(\pi\pr w\) the profit of a firm paying wage $w$. Naturally,
    \(\Pi\pr{w\pr t} = \pi\pr{w\pr t+t}+\pr{n-1}\pi\pr{w\pr t}\), and
    \[\ba
    & \frac{\partial \Pi\pr{w\pr t}}{\partial t} = \pr{w'\pr t+1}\cdot\pi'\pr{w\pr t+t}+\pr{n-1}w'\pr t\cdot \pi'\pr{w\pr t}\\
    \too & \frac{\partial \Pi\pr{w\pr t}}{\partial t}\bigg|_{t\to 0} = \pr{w'\pr 0+1}\cdot \pi'\pr{w\pr 0}+\pr{n-1}w'\pr 0\cdot \pi'\pr{w\pr 0} = 0
    \ea\]
  \item
    From quantitative analysis we can see that, the firms as a whole do
    not bear the tax burden and the workers bear all the tax burden.
  \item
    Intuitively, it must be the workers that bear all the tax burden.
    Even though the labor supply for any firm is perfectly elastic, the
    \textbf{total} labor supply is \textbf{perfectly inelastic}. Small as the impact on the wages in other towns, it is not
    negligible. Indeed, the labor supply for any given firm is perfectly
    elastic, so we cannot ignore the general equilibrium effect.
  \end{itemize}
\end{itemize}
}

\section{Pure Exchange Economy}\label{pure-exchange-economy}

A full market analysis typically involves three activities:
consumption, production, and trade. The cleanest setting to develop
general-equilibrium intuition strips out production: in a \emph{pure
exchange economy}, each agent starts with an endowment of goods, and
the only economic activity is mutually beneficial trade.

\ex{
Consider a perfectly competitive economy with two agents \((i = A, B)\)
and two goods \((j = 1, 2)\). Suppose the agents' preference relations
and initial endowments are given by
\[\ba
& u^A\pr{x_1,x_2} = x_1^2x_2\text{ with }e^A = \pr{1,2}\\
& u^B\pr{x_1,x_2} = x_1x_2^2\text{ with }e^B = \pr{2,1}
\ea\]

Can we come up with a price vector \(\pr{p_1,p_2}\) that clear both
markets?
}

\sol{
The key idea is that, each agent has their own income, which is endogenously given by the equilibrium price vector.

Suppose there is an equilibrium price vector \(\pr{p_1,p_2}\), then
agent \(A\)'s ``income" is \(m^A = p_1+2p_2\) and agent
\(B\)'s ``income" is \(m^B = 2p_1+p_2\). From this we can
pin down the optimal individual demands
\[\ba
\x^A &= \pr{\frac{2\pr{p_1+2p_2}}{3p_1},\frac{1\pr{p_1+2p_2}}{3p_2}}\\
\x^B &= \pr{\frac{1\pr{2p_1+p_2}}{3p_1},\frac{2\pr{2p_1+p_2}}{3p_2}}
\ea\]

In equilibrium, market demand equals market endowment for either good, respectively:
\[\ba
\frac{2\pr{p_1+2p_2}}{3p_1} + \frac{1\pr{2p_1+p_2}}{3p_1} & = 3\\
\frac{1\pr{p_1+2p_2}}{3p_2} + \frac{2\pr{2p_1+p_2}}{3p_2}& = 3
\ea\]

From either equation, we have \(\dfrac{p_1^*}{p_2^*} = 1\). (Otherwise
if there comes with a contradiction, the equilibrium cannot exist.)
}


From this example, we obtain four observations in this two-agent, two-good pure exchange economy:

\begin{enumerate}
\item
  There is an equilibrium price vector that clears both markets.
\item
  Only the relative price matters in equilibrium.
\item
  When one market clears, the other market clears simultaneously (and miraculously).
\item
  The equilibrium allocation is Pareto efficient.
\end{enumerate}

\subsection{Setups}\label{basic-setups-2}

\textbf{Model Preliminaries}

\begin{itemize}
\item
  \(I = \brace{1,2,\cdots,n}\) agents and \(J = \brace{1,2,\cdots,m}\)
  goods.
\item
  \(\pref^{i}\): preference relation of agent \(i\).
\item
  \(\textbf e^i = \brace{e_1^i,e_2^i,\cdots,e_m^i}\) initial endowment
  of agent \(i\) (property rights are well-defined and no co-ownership).
\item
  Denote the economy
  \(\mathcal E= \brace{\pref^i,\textbf e^i}_{I\in \mathcal I}\).
\end{itemize}

\asm{Market Structure}{
\begin{enumerate}
\item
  Perfect and complete information.
\item
  Perfectly competitive markets.

  \begin{itemize}
  \item
    Agents are price-takers.
  \item
    Prices are linear.
  \end{itemize}
\item
  Goods are perfectly divisible.
\end{enumerate}
}

\asm{Agents' Preferences and Endowments}{
For any agent \(i\in \mathcal I\):
\begin{enumerate}
\item
  The preference relation \(\pref^i\) is rational (complete and
  transitive) and continuous.
\item
The preference relation \(\pref^i\) is monotonic.
\item
   The preference relation \(\pref^i\) is (weakly) convex.
\item
   \(e_j^i>0\), for all \(j\in \mathcal J\).
\end{enumerate}
}

Notice that by the first assumption, the utility representation of $\pref^i$, say \(u^i\pr{\cdot}\), 
is guaranteed for any \(i \in I\), so we can alternatively denote the
economy as \(\mathcal E = \pr{u^i \pr\cdot,\e^i}_{i\in\mathcal I}\).

\subsection{Walrasian Equilibrium}\label{walrasian-equilibrium}

\defn{Walrasian Equilibrium}{
A \textit{Walrasian
Equilibrium} for a perfectly competitive pure exchange economy
\(\mathcal E = \pr{u^i\pr{\cdot},\textbf e^i}_{i\in \mathcal I}\) is a
price vector \(\p\ge \0\) and an allocation
\(\pr{\x^i}_{i\in \mathcal I}\) such that:
\begin{enumerate}
\item
  Utility maximization: Agent \(i\in \mathcal I\) solves:
  \[\max_{\x^i\ge\0}u^i\pr{\x^i} \text{ s.t. }\p\cdot \x^i\le \p\cdot \textbf e^i\]
\item
  Market clearing for all goods:
  \[\sum_{i\in \mathcal I} \x^i\pr{\p} = \sum_{i\in \mathcal I}\textbf e^i\]
\end{enumerate}
}

Under Walrasian equilibrium, all markets clear at the same time. A
Walrasian equilibrium specifies both \textit{equilibrium price vector} and \textit{equilibrium allocation}.

For simplicity, suppose for any \(i\in \mathcal I\), \(\x^i\pr{\p}\) is
always unique. This is guaranteed if we assume each
\(\pref_i\) is strictly convex and \(\p\gg \0\).

\defn{Excess Demand}{
The (aggregate) \textit{excess demand} for good \(j\) is given by
\[z_j\pr{\p} = \sum_{i\in \mathcal I} x_j^i\pr{\p}-\sum_{i\in\mathcal I}e_j^i\]
}

With this formulation, a Walrasian equilibrium is achieved when each agent maximize their utility subject to budget constraint prescribed by their endowment, and
\[z_j\pr{\p} = 0,\forall j\in \mathcal J\]

\propp{Properties of Excess Demand}{
Let
\(\mathcal E\) be a perfectly competitive pure exchange economy with each agent's preference relation being rational, continuous and monotonic, then the aggregate excess demand function
\(\z\pr{\p}\) satisfies:
\begin{enumerate}
\item
  \textbf{Homogeneity of degree \(0\)}: \(z_j\pr{\p}\) is homogeneous of degree
  \(0\) in \(\p\), for any \(j\in \mathcal J\).
\item
  \textbf{Walras' law}: \(\p\cdot\z\pr{\p} = 0\).
\end{enumerate}
}{
\begin{itemize}
\item
  Homogeneous of degree \(0\)
  \begin{itemize}
  \item
    \(\x^i\pr\p\) solves
    \(\max_{\x^i\ge\0}u^i\pr{\x^i}\text{ s.t. }\p\cdot \x^i\le \p\cdot \e^i\).
  \item
    For any \(t>0\), the solution must be the same as:
    \(\max_{\x^i\ge\0}u^i\pr{\x^i}\text{ s.t. }\pr{t\p}\cdot \x^i\le \pr{t\p}\cdot \e^i\).
    Hence, \(\x^i\pr{\p} = \x^i\pr{t\p}\).
  \item
    For any \(t>0\),
    \(z_j\pr{t\p} = \sum_{i = 1}^n x_j^i\pr{t\p}-\sum_{i = 1}^n e_j^i =  \sum_{i = 1}^n x_j^i\pr{\p}-\sum_{i = 1}^n e_j^i = z_j\pr\p\).
  \end{itemize}
\item
  Walras' law
  \begin{itemize}
  \item
    Budget constraint must be binding for any agent \(i\in \mathcal I\),
    \(\p\cdot \x^i\pr\p = \p\cdot \e^i\). (This must hold because of assumption of 
    monotonicity of preference relation, which is stronger than locally
    non-satiation).
  \item
    It follows that
    \[\ba
    \p\cdot \z\pr\p & = \p\cdot \pr{\sum_{i = 1}^n \x^i\pr\p - \sum_{i = 1}^n \e^i}\\
    & = \sum_{i = 1}^n \pr{\p\cdot \x^i\pr\p-\p\cdot \e^i}\\
    & = 0
    \ea\]
  \end{itemize}
\end{itemize}
}

The following proposition, the existence of Walrasian equilibrium, is characterized as the cornerstone of modern economics.

\propp{Existence of Walrasian Equilibrium}{
Let
\(\mathcal E = \pr{u^i\pr\cdot,\e^i}_{i\in\mathcal I}\) be a perfectly
competitive pure exchange economy that satisfies all the four assumptions on agents' preferences and endowments. Then
a Walrasian equilibrium exists.
}{
We give proof for the case of \(m = 2\) and \(\pref^i\)
being strictly convex for any \(i\in\mathcal I\). Given rationality, continuity and 
strict convexity of \(\pref^i\), \(\x^i\pr{p_1,p_2}\) is always unique
for any \(\p\gg\0\). When \(m = 2\), Walras' law states 
that
\[p_1z_1\pr{p_1,p_2}+p_2z_2\pr{p_1,p_2} = 0.\]
Then a Walrasian equilibrium is given by \(z_1\pr{p_1,p_2} = 0\) or
\(z_2\pr{p_1,p_2} = 0\). By monotonicity and strictly convexity of preference relation, the only possibility is
\(p_1,p_2>0\) in equilibrium, so we can normalize \(p_1 = 1\) or
\(p_2 = 1\).

By theorem of the maximum, given rationality, continuity, monotonicity and strict convexity of preference relation, \(x_i^1\pr{p_1,1}\) is
continuous for any \(i\in \mathcal I\). Hence, \(z_1\pr{p_1,1}\) is
continuous on \(\pr{0,+\infty}\). Given monotonicity and strict
convexity, \(u^i\pr\cdot\) is strongly monotone. Therefore,
\(\lim\limits_{p_1\to 0+}x_1^i\pr{p_1,1} = +\infty\) and
\(\lim\limits_{p_1\to0+}z_1\pr{p_1,1} = +\infty\). Symmetrically, we
have \(\lim\limits_{p_2\to 0+}x_2^i\pr{1,p_2} = +\infty\) and
\(\lim\limits_{p_2\to0+}z_2\pr{1,p_2} = +\infty\). By continuity,
\(\exists 0<\underline p<\bar p\) such that
\(z_1\pr{\underline p,1}>0>z_1\pr{\bar p,1}\). By the intermediate value
theorem, \(\exists p_1^*\in \pr{\underline p,\bar p}\) such that
\(z_1\pr{p_1^*,1} = 0\).
}

\section{Allocation}\label{allocation}

Recall the two cornerstone results from intermediate micro:
\begin{itemize}
\item
  \emph{First Welfare Theorem.} Any Walrasian equilibrium allocation
  is Pareto efficient.
\item
  \emph{Second Welfare Theorem.} Any Pareto efficient allocation can
  be supported as a Walrasian equilibrium after a suitable
  redistribution of initial endowments.
\end{itemize}

We now state these rigorously and strengthen the first. Begin with
the formal definition of a feasible allocation.

\defn{Feasible Allocation}{
For a pure exchange economy
\(\mathcal E\), an allocation \(\pr{\x^i}_{i = 1}^n\) with $\x^i\ge\0$ for all $i$ is
\textit{feasible} if for any good \(j\),
\[\sum_{i = 1}^n x_j^i\le \sum_{i=1}^n e_j^i.\]
}

In other words, total consumption of every good cannot exceed the
total endowment of that good.

\defn{Pareto Efficient Allocation}{
Given an economy
\(\mathcal E\), a feasible allocation \(\pr{\x^i}_{i = 1}^n\) is
(strongly) \textit{Pareto efficient} if there is no other feasible
allocation \(\pr{\w^i}_{i = 1}^n\) such that \(\w^i\pref^i \x^i\) for
all \(i\in \mathcal I\), with \(\w^k\succ^k \x^k\) for some
\(k\in \mathcal I\).
}

A Pareto efficient allocation is one from which no agent can be made
strictly better off without making at least one other agent strictly
worse off.

\defn{Core Allocation}{
A feasible allocation
\(\pr{\x^i}_{i = 1}^n\) is \textit{in the core} of \(\mathcal E\) if
there is no group \(\mathcal I_0\se \mathcal I\) and an alternative allocation
\(\pr{\w^i}_{i = 1}^n\) such that:
\begin{enumerate}
\item
  For any \(j\in \mathcal J\),
  \(\sum_{i\in \mathcal I_0}w_j^i\le \sum_{i\in \mathcal I_0}e_j^i\).
\item
  \(\w^i\pref^i\x^i\) for all \(i\in \mathcal I_0\), with
  \(\w^k\succ^k\x^k\) for some \(k\in \mathcal I_0\).
\end{enumerate}
}

An allocation is in the core if \emph{no} group of agents — large or
small — can find a redistribution of their own endowments that makes
every member of the group at least as well off and some member strictly
better off. Core allocations are Pareto efficient (taking the
``grand coalition'' as the blocking group recovers the Pareto
condition), but Pareto efficiency does not imply the core: an
allocation can be Pareto efficient overall yet have a strictly
sub-coalition that could improve on it among themselves.

\thmp{Strengthened First Theorem of Welfare Economics}{
Let \(\mathcal E\) be a perfectly competitive pure exchange
economy in which each agent's preference relation is rational, continuous and concave. Denote \(S_P\) as the set of Pareto efficient allocations,  \(S_C\) the set of core allocations, and \(S_W\) the set of Walrasian equilibrium allocations. Then
\[S_W\subseteq S_C\subseteq S_P\]
}{
By definition, \(S_C\subseteq S_P\), so it
suffices to show \(S_W\subseteq S_C\). Suppose not, then there exists a
Walrasian equilibrium \(\pr{\p;\pr{\x^i}_{i\in \mathcal I}}\), a group
\(\mathcal I_0\se \mathcal I\) and an alternative allocation
\(\pr{\w^i}_{i\in\mathcal I}\) such that:
\[\bc
\w^i\pref^i \x^i,\text{ for all }i\in \mathcal I_0\\
\w^k\succ^k \x^k, \text{ for some }k\in \mathcal K_0\subseteq \mathcal I_0
\ec\]
By direct revealed preference,
\[\bc
\p\cdot \w^k> \p\cdot \x^k,\text{ for all }k\in \mathcal K_0\\
\p\cdot \w^i\ge \p\cdot \x^i,\text{ for all }i \in \mathcal I_0\backslash \mathcal K_0
\ec\]

Summing over \(i\in \mathcal I_0\), we have
\[\ba
& \p\cdot \sum_{i\in \mathcal I_0}\w^i> \p\cdot \sum_{i\in \mathcal I_0}\x^i = \p\cdot \sum_{i\in\mathcal I_0}\e^i\\
\too & \sum_{i\in\mathcal I_0}\w^i\nleq \sum_{i\in\mathcal I_0}\e^i
\ea\]
}

\rmkb{
\begin{itemize}
\item
  The standard First Welfare Theorem gives \(S_W\subseteq S_P\); the
  strengthened version shows the tighter \(S_W\subseteq S_C\).
\item
  \(S_W\subseteq S_C\) says that any Walrasian equilibrium is not only
  efficient but also \emph{coalition-proof}: no group can do better
  for itself by trading internally. This is a stronger fairness
  justification for the price mechanism than mere Pareto efficiency.
\item
  The theorem demands very little of preferences (rationality,
  continuity, concavity). What it does need is the institutional
  scaffolding: perfect competition, complete information, no
  externalities, complete markets. Drop any of these and the
  conclusion can fail.
\end{itemize}
}

\thm{Second Theorem of Welfare Economics}{
Let
\(\mathcal E\) be a perfectly competitive pure exchange economy in which each agent's preference relation is rational, continuous and concave. If \(\x^i\gg\0\) is Pareto
efficient, for all \(i\), then there exists an initial endowment
\(\pr{\e^i}_{i \in\mathcal I}\) and a price vector \(\p\ge\0\) such that
\(\pr{\p;\pr{\x^i}_{i\in \mathcal I}}\) is a Walrasian equilibrium given
these endowments.
}

The Second Welfare Theorem is, in a sense, less operational than the
First. To support a given Pareto-efficient allocation, the planner
needs both the political authority to redistribute initial endowments
freely \emph{and} full information about every agent's preferences —
neither of which is generally available in practice.

\section{General Equilibrium with Production}\label{ge-with-production}

\subsection{Setups}

Recall that in producer theory, we used production sets and ownership
shares to describe the firms and their production technologies.

Suppose there are \(K\) firms (\(\mathcal K = \brace{1,2,\cdots,K}\)) in the economy, each firm $k\in \mathcal K$ with its production set \(Y^k\). Let \(\alpha^{ki}\ge0\) be agent
\(i\)'s ownership share of firm \(k\). The production
economy can then be described as
\[\mathcal E = \pr{\pr{u^i,\e^i,\pr{\alpha^{ki}}_{k\in \mathcal K}}_{i\in\mathcal I},\pr{Y^k}_{k\in \mathcal K}}.\]

\defn{Walrasian Equilibrium with Production}{
A
\textbf{Walrasian Equilibrium} for a perfectly competitive production
economy
\(\mathcal E = \pr{\pr{u^i,\e^i,\pr{\alpha^{ki}}_{k\in \mathcal K}}_{i\in\mathcal I},\pr{Y^k}_{k\in \mathcal K}}\)
is a vector
\(\pr{\p,\pr{\x^i}_{i\in \mathcal I},\pr{\y^k}_{k\in\mathcal K}}\) such
that:

\begin{enumerate}
\item
  \textbf{Profit maximization}: Firm \(k\) solves
  \[\max_{y^k\in Y^k}\p\cdot \y^k\]
\item
  \textbf{Utility maximization}: Agent \(i\) solves
  \[\max_{\x^i\ge \0}u^i\pr{\x^i} \text{ s.t. }\p\cdot \x^i\le \p\cdot \e^i+\sum_{k\in \mathcal K}\alpha^{ik}\p\cdot \y^k\]
\item
  \textbf{Market clearing for all goods}:
  \[\sum_{i\in \mathcal I}\x^i\pr \p = \sum_{i\in \mathcal I}\e^i+\sum_{k\in \mathcal K}\y^k\pr\p\]
\end{enumerate}
}

That is, a Walrasian equilibrium specifies prices, consumption
bundles, and production plans such that every agent maximizes utility
(taking firm profits as given dividend income), every firm maximizes
profit, and all markets clear simultaneously.

The conceptual structure is identical to the pure-exchange case;
production simply adds the firms' optimization problems. For the
existence theorem to go through, however, we need three regularity
conditions on production technologies.

\asm{Production Technology}{
For each firm \(k\in \mathcal K\),
\begin{enumerate}
\item
  \(Y^k\ne\varnothing\) is \textbf{closed
  and convex}.
\item
   \textbf{Shutdown and free disposal}: that is, \(\0\in Y^k\), and \(\y\in Y^k\) implies
  \(\y'\in Y^k\), for all \(\y'\le \y\).
\item
  \textbf{Irreversibility}: Let \(Y = \bigcup_{k\in \mathcal K} Y^k\), then
  \(Y\cap \pr{-Y} = \brace{\0}\).
\end{enumerate}
}

\rmkb{
\begin{itemize}
\item
  \(Y^k\ne\varnothing\) is innocuous, otherwise since firm \(k\) is allowed to produce
  nothing, we can just discard \(Y^k\) if $Y^k = \varnothing$.
\item
  Irreversibility says that the production cannot be completely
  reversed.
\end{itemize}
}

\subsection{Existence of Walrasian Equilibrium with Production}\label{existence-of-walrasian-equilibrium-with-production}

\thm{Existence of Walrasian Equilibrium with Production}{
Let
\(\mathcal E = \pr{\pr{u^i,\e^i,\pr{\alpha^{ki}}_{k\in \mathcal K}}_{i\in\mathcal I},\pr{Y^k}_{k\in \mathcal K}}\)
be a perfectly competitive production economy satisfying all the assumptions on preference relation and endowments, and those on  production technology. Then a Walrasian equilibrium
\(\pr{\p,\pr{\x^i}_{i\in\mathcal I},\pr{\y^k}_{k\in\mathcal K}}\)
exists.
}

\rmkb{
\begin{itemize}
\item
  This proposition of existence is the extension of the existence result on pure exchange economies.
\item
  The first and second theorem of welfare economics also generalize to
  economics with production.
\end{itemize}
}

Recall that when a production technology exhibits constant returns to scale,
the firm's profit must be either \(0\) or \(+\infty\). For
market clearing to make sense, it must be that each firm is making \(0\) profit.

For simplicity, consider an economy with two goods. The production technology can  linearly transform \(b\) units of good 1 into
(at most) \(c\) units of good 2. The
production set can then be depicted as
\[Y = \brace{\pr{y_1,y_2}|y_1\le 0,y_2\ge 0,y_2\le -\frac{c}{b}y_1}.\]
Alternatively, we can present the production technology with the
vector \(a = \pr{-1,\dfrac{c}{b}}\).

If a linear activity production $a = (a_1,a_2)$ is ever used in a Walrasian equilibrium
(i.e., \(\l^*>0\)), then it must be that (because the firm is making zero profit)
\[\pr{p_1^*,p_2^*}\cdot \pr{a_1,a_2} = 0.\]

\ex{
Consider a perfectly competitive economy \(\mathcal E\)
with two goods (1 and 2) and two agents (\(A\) and \(B\)). The
agents' utility functions and initial endowments are as
follows:
$$
\begin{aligned}
 u^A\pr{x_1,x_2} = x_1x_2,&\  \e^A = \pr{1,0}\\
 u^B\pr{x_1,x_2} = x_1x_2^2, &\  \e^B=\pr{0,1}
\end{aligned}
$$

First suppose the economy is pure exchange with no production.

\begin{enumerate}
\item
  Derive the set of Pareto efficient allocations.
\item
  Derive the set of core allocations.
\item
  Derive a Walrasian equilibrium
  \(\pr{\pr{p_1,p_2};\pr{x_1^A,x_2^A,x_1^B,x_2^B}}\).
\end{enumerate}

Next we introduce production: Suppose there is a perfectly competitive
firm which can transform one unit of good 1 into (at most) one unit of
good 2, i.e., \(\textbf a =\pr{-1,1}\).

\begin{enumerate}
\item[4.]
  Derive a Walrasian equilibrium with production
  \(\pr{\pr{p_1,p_2,};\pr{x_1^A,x_2^A,x_2^B,x_2^B},\lambda}\).
\end{enumerate}
}

\sol{
\begin{enumerate}
\item Since both agents' preference
relations are monotonic, any Pareto efficient allocation must satisfy
\[\x^A+\x^B = \e^A+\e^B\]
Clearly, \(\pr{\pr{1,1},\pr{0,0}}\) and \(\pr{\pr{0,0},\pr{1,1}}\) are
Pareto efficient.

\rmk{
\(\pr{\pr{1,0},\pr{0,1}}\) or
\(\pr{\pr{0,1},\pr{1,0}}\) are not necessarily Pareto efficient, though we cannot
make any agent better off material-wise without making the other worse off material-wise. However, agents' well-being is measured in terms of magnitude of utilities.
}

Apart from corner solutions, by the idea of no gain from trade in equilibrium, we must have
\[|\mathit{MRS}^A| = |\mathit{MRS}^B|\]
Therefore, the set of efficient allocations is given by
\[S_P = \brace{\pr{x_1^A,x_2^A,x_1^B,x_2^B}\ge \0: x_1^Ax_2^B = 2x_2^Ax_1^B,x_1^A+x_1^B = x_2^A+x_2^B = 1}.\]
It has been checked that \(\pr{\pr{1,1},\pr{0,0}}\) and \(\pr{\pr{0,0},\pr{1,1}}\) are included in the set $S_P$.
\item When there are only two agents, the only subgroups other than the grand
coalition are \(\brace 1\) and \(\brace 2\). In other words, on the basis of Pareto efficiency, the additional requirement for an allocation to be in the core is that {no agent is made worse in any allocation compared to their own endowment}, which is also termed \textit{individual
rationality}. Since the initial endowment is the worst possible
allocation for each agent, the set of core allocations \(S_C = S_P\).
\item Suppose there is an
equilibrium price vector \(\pr{p_1,p_2}\). Then agent
A's ``income" \(m^A = p_1\) and agent B's
``income" \(m^B = p_2\). The optimal individual demands are
\[\ba
\pr{x_1^A,x_2^A} & = \pr{\frac12,\frac{p_1}{2p_2}}\\
\pr{x_1^B,x_2^B} & = \pr{\frac{p_2}{3p_1},\frac23}
\ea\]
In equilibrium, market demand equals total endowments for each good.
\[\bc
\frac12+\frac{p_2}{3p_1} = 1\\
\frac{p_1}{2p_2}+\frac{2}{3} = 1
\ec\too \frac{p_1^*}{p_2^*} = \frac23\]
Hence, a Walrasian equilibrium is
\(\pr{\p = \pr{2,3},\x^A = \pr{\frac{1}{2},\frac13},\x^B = \pr{\frac12,\frac23}}\).

\rmk{
Here ``\textbf{a} Walrasian equilibrium" is emphasized because only relative prices matter in Walrasian equilibrium; theoretically there could be infinitely-many equilibria in terms of absolute prices.
}
\item Suppose the production technology is used in equilibrium, then we must
have \(\pr{p_1,p_2}\cdot \textbf a = \0\), that is, \(p_1 = p_2\iff \frac{p_1}{p_2} = 1\).

In equilibrium, market demand equals market supply for each good.
\[\bc
\frac12+\frac13 = 1-\l\\
\frac12+\frac23 = 1+\l
\ec\too \l^* = \frac16\]
A Walrasian equilibrium is given by
\(\pr{\p^* = \pr{1,1},\x^A = \pr{\frac12,\frac12},\x^B = \pr{\frac13,\frac23},\l^* = \frac16}\).
\end{enumerate}
}

